\section{The sharp error for every cyclic cut of size at most six}
\label{sec:root-all-six-cuts}

The bounded-shore case needs an exact counting deficit.
A missing forest in one parity class supplies that deficit.

\begin{proposition}[A bounded six-pole shore of girth at least seven]
\label{prop:root-bounded-six-shore}
Fix an integer $H\ge1$.
Let normalized six-poles $X,Y$ have ranks $1\le h\le H$ and $k$, respectively.
Assume that $X$ has girth at least seven.
Join corresponding ports to form a cubic graph $G$ of rank $r=h+k+5$.
Then, uniformly over these choices, for all sufficiently large $k$,
\[
 c(G)<2^{r-1}.
\]
\end{proposition}

\begin{proof}
A planar normalized six-pole cannot have girth at least seven, by the Euler argument in Proposition~\ref{prop:root-six-girth-path-gap}.
Thus $X$ is nonplanar.
Theorem~\ref{thm:path-six-full-four-planar} supplies a four-port set $S_0$ for which $F_{S_0}(X)<2^h$.
Write $\delta=2^{-h}$.
Integrality gives
\[
 2^{-h}F_{S_0}(X)\le1-\delta.
\]
No terminal pair in $X$ is saturated, by Lemma~\ref{lem:root-saturated-pole-planar}.
Integrality again gives $2^{-h}P_e(X)\le1-\delta$ for every terminal pair $e$.
All other normalized forest totals are at most one.

Use $a_e,U_S,W$ for the normalized path, four-terminal, and six-terminal totals of $Y$.
Let
\[
 A_Y=\sum_e a_e+\sum_SU_S+W,
 \qquad T_Y=\sum_{S\ne S_0}U_S+W.
\]
There are fourteen four-port terms in $T_Y$, and one six-port term.
Consequently $T_Y\le15$.
The fixed-port aggregate theorem gives $A_Y\le16+\eta_k$, with $\eta_k\to0$ uniformly.

Drop the incompatible-pairing subtractions in the exact six-cut crossing expression.
The coefficients of all path sectors and of $S_0$ are at most $1-\delta$.
Every remaining coefficient is at most one.
Therefore the normalized crossing count satisfies
\[
 \mathcal B(X,Y)
 \le(1-\delta)A_Y+\delta T_Y
 \le16-\delta+(1-\delta)\eta_k.
\]
This estimate uses the exact quantum $\delta$, rather than an unspecified strict inequality.

The asymptotic cycle bound gives $c(Y)\le(1/2+\theta_k)2^k$, where $\theta_k\to0$ uniformly.
The trivial bound gives $c(X)\le2^h-1$.
Thus
\[
 \frac{c(G)}{2^{h+k}}
 \le16-\frac\delta2+(1-\delta)\eta_k
       +\delta\theta_k+2^{-k}.
\]
Here $\delta\ge2^{-H}$.
Choose $k$ sufficiently large that the three error terms total less than $2^{-H-1}$.
The displayed quantity is then strictly below sixteen, uniformly in $X$ and $h\le H$.
\end{proof}

\begin{theorem}[A strict deficit for genuine six-cuts]
\label{thm:root-six-cut-strict}
There is an absolute integer $R_6$ with the following property.
Let $G$ be a connected loopless cubic multigraph of rank $r\ge R_6$.
Assume that $G$ has a cyclic cut of size six and no cyclic cut of size at most five.
Then
\[
 c(G)<2^{r-1}.
\]
\end{theorem}

\begin{proof}
The cut shores are connected normalized six-poles, as proved in Theorem~\ref{thm:four-six-two-large-shores}.
Write their ranks as $h\le k$.

First suppose that $G$ has a cycle of length at most six.
For sufficiently large rank, its complement contains a cycle.
A shortest such cycle is induced.
A cycle of length at most five would therefore give a forbidden smaller cyclic cut.
If its length is six, the six-cycle-shore theorem gives the asserted strict deficit for sufficiently large rank.
The normalization of the opposite shore follows from the absence of smaller cyclic cuts.

We may therefore assume that $G$ has girth at least seven.
Let $H$ be the threshold from Theorem~\ref{thm:four-six-two-large-shores}.
If $h,k\ge H$, that theorem proves the claim.
Otherwise $h<H$ and $k=r-h-5$ tends to infinity with $r$.
Proposition~\ref{prop:root-bounded-six-shore}, with the fixed bound $H$, proves the claim uniformly.
Increase $R_6$ to meet all these thresholds.
\end{proof}

\begin{corollary}[Sharp order for cuts of size at most six]
\label{cor:root-all-six-cuts-sharp}
There is an absolute constant $K$ such that every connected loopless cubic multigraph of rank $r$ with a cyclic cut of size at most six satisfies
\[
 c(G)\le2^{r-1}+K r2^{r/2}.
\]
The order of the error is necessary within this class, by the Horak--Kim cyclic-square family.
\end{corollary}

\begin{proof}
For a cyclic cut of size at most five, apply the established sharp small-cut theorem.
For a six-cut without any smaller cyclic cut, apply Theorem~\ref{thm:root-six-cut-strict} when $r\ge R_6$.
The finitely many remaining ranks are covered by increasing $K$, using $c(G)\le2^r-1$.
The cited family has cyclic four-cuts and excess of order $r2^{r/2}$.
\end{proof}

The unrestricted sharp bound is still open.
Its remaining large cubic kernels may now be assumed to have no cyclic cut of size at most six and to have girth at least seven.
